\section[From business phrases to causal variables]{From business phrases to causal variables and defensible diagrams}
\label{sec:exec28-cm}

A useful causal assistant must decide what the variables mean before it can reason about arrows. In a servicing discussion, current borrower, delinquent borrower and financially distressed borrower may denote values of different operational definitions, overlapping classifications or measurements taken at different times. Automatically drawing one node for each phrase can turn a vocabulary mistake into a scientific model. The first task is to propose a coherent representation and preserve alternatives that cannot yet be resolved.

\subsection{Two concrete language tasks}

The inspected preprint of \citet{exec28_cm_razouk} evaluates two tasks directly relevant to this problem. The first asks whether two textual entities represent different values of the same variable, without supplying that variable's name. The second asks whether an entity refers simultaneously to values of several supplied variables. The source calls the second object an interaction entity. These are useful semantic tasks, distinct from verifying a causal arrow, identifying an effect or establishing an intervention's feasibility.

For a bank, consider whether payment within the grace period and payment after the grace period are values of a common timing variable. The answer depends on the account, payment obligation, time origin and grace-period definition. A phrase describing a mortgage customer with a payment holiday combines product and servicing status. Detecting that combination does not establish that these statuses interact in their causal effect on default. The assistant should first return candidate variable definitions, value domains, unit identifiers, observation clocks and textual evidence. A subsequent model specifies relationships among those variables.

A proposed representation can be written as
\begin{equation}
 \mathcal V_j=(\text{unit},\text{quantity},\text{domain},\text{time},
                \text{measurement},\text{availability}),
\end{equation}
where a phrase may map to a value, a predicate on a domain or an unresolved collection of variables. This is a product design convention, not a theorem that language has a unique causal decomposition. A high-risk label can denote a score category, a decision-maker's belief or an unobserved propensity. Treating them as interchangeable can make a score's prediction appear to be an intervention on underlying risk.

Pairwise language answers also need consistency review. Suppose an assistant groups phrase $e_0$ with $e_1$, groups $e_1$ with $e_2$, and rejects grouping $e_0$ with $e_2$. Those answers cannot all define an equivalence relation for a single partition. The proposed product should expose the conflicting triangle and relevant definitions. Taking the transitive closure silently can merge incompatible variables, while deleting one pair without a reason hides the conflict. Multiple candidate representations are preferable when the language remains ambiguous. The independent finite reference reproduces the contradiction.\citep{exec28_cm_reference}

\subsection{What the evaluation actually establishes}

The source uses a within-subjects design in which seven named language models generate examples and evaluate examples from themselves and other models. It sets temperature to zero and uses the same English prompts. Its domain tables contain eight domains, including finance. Task 1 uses 560 positive and 560 negative examples; Task 2 uses 559 positive and 560 negative examples. These are reported evaluation denominators, not independently audited bank expert labels.\citep{exec28_cm_razouk}

The results illustrate why a single model ranking is insufficient. In the source's Task 1 table, GPT-4-turbo has reported precision 97\%, recall 41\% and kappa 40\%; rankings differ for the second task. These historical labels and reported numbers describe that generated-data evaluation. They are not a recommendation for a current provider, a service guarantee or evidence that 97\% of suggested banking variables are correct. High precision with low recall can leave important distinctions unresolved. The cost of an omitted variable may differ greatly from that of a false proposed merge.\citep{exec28_cm_razouk}

The paper explicitly states that its numbers derive from generated data and need comparison with domain-specific expertise. Agreement with generation-derived labels can reward shared conventions or mistakes. Cross-model evaluation broadens the comparison but does not produce independently established causal truth. Temperature zero supplies neither calibrated uncertainty nor an immutable service implementation. Domain dependence is a reason to stratify later assessment by banking journey, variable type, language and ambiguity rather than transport an overall score to every workflow.

A proposal-stage evaluation can retain inexpensive generated fixtures alongside a small, independently adjudicated public set. Generated fixtures make intended variable/value relationships explicit and support controlled perturbations. Public documents provide realistic ambiguity, missing definitions and conflicting timestamps. Their reference labels need documented provenance; disagreement should be preserved instead of resolved by asking the same model to grade itself. The output includes an unresolved option, with performance reported separately for accepted, rejected and deferred cases. Coverage, incorrect merges and splits, unsupported certainty, human review effort and full workflow cost all matter.

\subsection{A compound phrase does not establish causal interaction}

Let $A$ denote offering a payment holiday and $B$ an outreach action. A phrase mentioning both actions does not determine their additive interaction in a monetary outcome. That requires a four-regime contrast,
\begin{equation}
 I=E[Y(1,1)]-E[Y(1,0)]-E[Y(0,1)]+E[Y(0,0)].
 \label{eq:exec28-cm-interaction}
\end{equation}
In an authored example, set $Y(a,b)=100+10a+20b$ dollars per account. The outcomes are 100, 110, 120 and 130, giving $I=0$ despite the compound phrase. Changing the stipulated law to $100+10a+20b+5ab$ changes the joint cell to 135 and gives $I=5$ dollars. Exact enumeration checks both contrasts. This is a scale-specific potential-outcome calculation, not measured servicing interaction.\citep{exec28_cm_reference}

A deterministic compound node $C=AB$ may be convenient, but the analyst must state what an intervention on $C$ means. Setting $C=1$ without specifying $A$ and $B$, or changing $C$ while holding its defining parents fixed, need not describe a feasible action. Multiple versions of a compound label can have different consequences. The assistant should require an intervention definition rather than infer it from a noun phrase.

The source invokes causal transitivity as motivation. That terminology should not become a numerical rule that every directed path has a nonzero total effect. Let $U$ be uniform on $\{0,1,2,3\}$, define $B=\mathbf 1\{U<1+2A\}$ and set $Y=100+20B-10A$. The model has an $A\to B\to Y$ path and a direct $A\to Y$ path. Under interventions on $A$, the mediator probability changes from $1/4$ to $3/4$. This produces a mediated mean change of USD~10 and a direct change of minus USD~10. Both intervention means are USD~105: the total mean effect is zero through cancellation, although both mechanisms are present.\citep{exec28_cm_reference}

This finite example distinguishes graph structure, mechanisms and a numerical contrast; it does not assert that the source's experiment is invalid. The proposed validator should check stated graph semantics, temporal order and intervention definitions, rather than impose a blanket effect-magnitude rule. Statistical dependence also cannot by itself endorse a causal arrow. A graph, an observation law and an estimand serve different purposes.

\subsection{An integration contract for semantic proposals}

The proposed CLI or MCP interface can expose variable construction as a reviewable artifact. Each variable should carry its source span, value domain, population, event clock, availability clock, measurement rule and possible action semantics. A mapping can be accepted, disputed or unresolved. Relations should distinguish source assertions, model-generated hypotheses, statistical associations and approved modeling assumptions. A syntactically valid DAG must not erase these distinctions.

The retained source links to an anonymous code-and-data repository. The link is an acquisition lead; it does not establish a pinned release, licence, model service identity or executable package in this workspace. Direct runtime adoption is deferred. The tasks remain useful leads for an independently specified prototype. Original prompts and reference labels are separate artifacts whose rights and provenance require review. A prototype using separately authorized fixtures can test representation quality without claiming to reproduce the source's runtime.

Before estimation, the assistant should present competing variable definitions, an intervention and outcome, unresolved timing or measurement issues, and the assumptions needed to compare selected policies. The business owner can then see what information is missing and why it matters. The result is a traceable scientific specification whose ambiguity remains visible, rather than a diagram that merely looks complete.
